\section{Описание практической части}
\label{sec:Chapter4} \index{Chapter4}

\subsection{Инструментарий и язык реализации}

Система реализована на языке \textbf{Python~3.10} с использованием фреймворка
\textbf{PyTorch~2.1}. Выбор обусловлен зрелостью экосистемы глубокого обучения,
поддержкой автоматического дифференцирования и доступностью предобученных
моделей. Применялись следующие библиотеки:
\begin{itemize}
    \item \texttt{torchaudio} --- загрузка аудио, ресемплинг, мел-преобразования;
    \item \texttt{transformers} (HuggingFace) --- предобученные веса WavLM Base+
        и HuBERT Base;
    \item \texttt{speechbrain} и официальная реализация HiFi-GAN --- вокодеры;
    \item \texttt{pesq}, \texttt{pystoi}, \texttt{torchmetrics} ---
        вычисление метрик качества;
    \item \texttt{numpy}, \texttt{scipy} --- генерация RIR и обработка сигналов.
\end{itemize}
Эксперименты проводились на одном GPU NVIDIA RTX~3090 (24~ГБ).

\subsection{Архитектура программного комплекса}

Программный комплекс организован по модульному принципу
(рис.~\ref{fig:codearch}). Модуль \texttt{data} формирует пары
<<чистый — искажённый>> на лету, применяя аугментации. Модуль \texttt{encoder}
инкапсулирует замороженный SSL-энкодер и извлечение признаков выбранного слоя.
Модуль \texttt{cleaner} реализует параллельные адаптеры~\eqref{eq:adapter}.
Модуль \texttt{vocoder} предоставляет единый интерфейс к HiFi-GAN и WaveFit.
Модуль \texttt{train} управляет циклом обучения, а \texttt{eval} ---
вычислением метрик.

\begin{figure}[H]
    \centering
    \begin{tikzpicture}[
        mod/.style={draw, rounded corners, minimum width=2.4cm,
            minimum height=0.8cm, align=center, font=\scriptsize, fill=black!6},
        >={Stealth[length=2.2mm]}, node distance=0.6cm and 0.7cm]
        \node (data) [mod] {\texttt{data}};
        \node (enc) [mod, right=of data] {\texttt{encoder}};
        \node (cln) [mod, right=of enc] {\texttt{cleaner}};
        \node (voc) [mod, right=of cln] {\texttt{vocoder}};
        \node (tr) [mod, below=of enc] {\texttt{train}};
        \node (ev) [mod, below=of cln] {\texttt{eval}};
        \draw[->] (data) -- (enc);
        \draw[->] (enc) -- (cln);
        \draw[->] (cln) -- (voc);
        \draw[->] (tr) -- (cln);
        \draw[->] (tr) -- (voc);
        \draw[->] (voc) |- (ev);
        \draw[->] (data) |- (tr);
    \end{tikzpicture}
    \caption{Модульная структура программного комплекса}
    \label{fig:codearch}
\end{figure}

\subsection{Процедура обучения}

Обучение проводилось в два этапа. На первом этапе при замороженном энкодере
обучались адаптеры очистителя по функции потерь~\eqref{eq:fcloss}. На втором
этапе вокодер дообучался на очищенных признаках с комбинированной функцией
потерь HiFi-GAN (состязательная, mel-spectrogram и feature matching). Основные
гиперпараметры приведены в табл.~\ref{tab:hparams}; полный перечень --- в
приложении.

\begin{table}[H]
    \caption{Основные гиперпараметры обучения}
    \label{tab:hparams}
    \centering
    \footnotesize
    \begin{tabularx}{\textwidth}{X l}
        \toprule
        Параметр & Значение \\
        \midrule
        Оптимизатор & AdamW ($\beta_1=0{,}8$, $\beta_2=0{,}99$) \\
        Скорость обучения & $2\cdot10^{-4}$ с экспоненциальным затуханием \\
        Размер батча & 16 \\
        Длительность сегмента & 2 с \\
        Число шагов (адаптеры) & 200 тыс. \\
        Число шагов (вокодер) & 400 тыс. \\
        Частота дискретизации & 16 кГц \\
        Скрытая размерность адаптера & 256 \\
        \bottomrule
    \end{tabularx}
\end{table}

Динамика обучения очистителя признаков для обоих энкодеров показана на
рис.~\ref{fig:trainloss}. Кривые демонстрируют монотонное убывание функции
потерь~\eqref{eq:fcloss} и выход на плато после $\sim$150~тыс. шагов; для WavLM
достигается стабильно меньшее значение потерь, что согласуется с его большей
устойчивостью к деградациям.

\begin{figure}[H]
    \centering
    \begin{tikzpicture}
    \begin{axis}[
        width=0.82\textwidth, height=5.8cm,
        xlabel={Шаг обучения, тыс.}, ylabel={Потеря $\mathcal{L}_{\text{fc}}$},
        xmin=0, xmax=200, ymin=0.1, ymax=0.9,
        legend pos=north east, legend cell align=left,
        grid=both, grid style={black!12},
        tick label style={font=\scriptsize}, label style={font=\small}]
    \addplot[thick] coordinates {
        (1,0.84)(10,0.56)(25,0.42)(50,0.31)(75,0.26)(100,0.22)(125,0.20)
        (150,0.19)(175,0.185)(200,0.182)};
    \addlegendentry{WavLM Base+}
    \addplot[thick, dashed] coordinates {
        (1,0.87)(10,0.61)(25,0.47)(50,0.36)(75,0.30)(100,0.27)(125,0.25)
        (150,0.24)(175,0.235)(200,0.232)};
    \addlegendentry{HuBERT Base}
    \end{axis}
    \end{tikzpicture}
    \caption{Кривые обучения очистителя признаков для двух энкодеров}
    \label{fig:trainloss}
\end{figure}

\subsection{Протокол экспериментов и воспроизводимость}

Для корректности сравнения все четыре конфигурации обучались и оценивались по
единому протоколу. Разбиение на обучающую, валидационную и тестовые выборки
(табл.~\ref{tab:data}) фиксировалось и не пересекалось по дикторам, что исключает
утечку информации между этапами. Генераторы случайных чисел инициализировались
фиксированным зерном, что обеспечивает повторяемость аугментаций при заданной
эпохе. Подбор гиперпараметров и номера слоя выполнялся исключительно на
валидационной выборке, а итоговые метрики сообщаются на тестовых наборах, ни разу
не использованных при настройке. Каждая метрика усреднялась по всем тестовым
примерам; для интрузивных метрик эталоном служил исходный чистый сигнал до
аугментации. Такой протокол гарантирует, что наблюдаемые различия между
конфигурациями обусловлены именно выбором энкодера, слоя и вокодера, а не
условиями эксперимента.

\subsection{Вычислительная сложность}

Основная вычислительная нагрузка приходится на прямой проход энкодера и вокодера.
Трансформерный энкодер имеет сложность $O(T^2 d)$ по длине последовательности
$T$ и размерности $d$ из-за механизма самовнимания, однако благодаря низкой
частоте кадров (50~Гц) и замороженным весам эта стоимость фиксирована и не
требует обратного распространения. Адаптеры добавляют лишь линейную по $T$
сложность $O(T d r)$ при размерности узкого места $r=256$. Вокодер HiFi-GAN
полностью свёрточный и имеет линейную по длине сигнала сложность; WaveFit
повторяет проход $N=5$ раз, что соответственно увеличивает время вывода.

\subsection{Результаты экспериментов}

\subsubsection{Сравнение конфигураций}

Четыре конфигурации (табл.~\ref{tab:configs}) обучены в идентичных условиях и
оценены на тестовых наборах. Усреднённые по тестовым выборкам результаты
приведены в табл.~\ref{tab:main}. Строка <<Вход (шум)>> соответствует
необработанному искажённому сигналу.

\begin{table}[H]
    \caption{Качество восстановления для исследуемых конфигураций
        (усреднение по тестовым наборам A и B)}
    \label{tab:main}
    \centering
    \footnotesize
    \begin{tabularx}{\textwidth}{l *{4}{>{\centering\arraybackslash}X}}
        \toprule
        Конфигурация & PESQ $\uparrow$ & STOI $\uparrow$ & SI-SDR, дБ $\uparrow$
            & DNSMOS $\uparrow$ \\
        \midrule
        Вход (шум)         & 1,82 & 0,841 & 4,7  & 2,41 \\
        \midrule
        K1 (HuBERT+HiFi-GAN) & 2,71 & 0,908 & 13,9 & 3,18 \\
        K2 (HuBERT+WaveFit)  & 2,78 & 0,912 & 14,3 & 3,29 \\
        K3 (WavLM+HiFi-GAN)  & \textbf{2,94} & \textbf{0,927} & 15,2 & 3,41 \\
        K4 (WavLM+WaveFit)   & 2,97 & 0,929 & \textbf{15,6} & \textbf{3,52} \\
        \bottomrule
    \end{tabularx}
\end{table}

Все конфигурации обеспечивают значительный прирост качества относительно входа:
PESQ возрастает на $0{,}9$--$1{,}15$ балла, STOI --- на $0{,}07$--$0{,}09$, что
превышает целевые ориентиры из раздела~\ref{sec:Chapter1}. Замена HuBERT на
WavLM даёт устойчивый прирост по всем метрикам (сравнение K1/K3 и K2/K4), что
подтверждает гипотезу о пользе денойзинга при предобучении энкодера. Переход от
HiFi-GAN к WaveFit улучшает метрики незначительно (на $0{,}03$ PESQ и $0{,}11$
DNSMOS в паре K3/K4).

\subsubsection{Вычислительная эффективность}

Время вывода и размер обучаемой части моделей приведены в
табл.~\ref{tab:efficiency}. Значение RTF измерено на RTX~3090 при обработке
10-секундных фрагментов.

\begin{table}[H]
    \caption{Вычислительная эффективность конфигураций}
    \label{tab:efficiency}
    \centering
    \footnotesize
    \begin{tabularx}{\textwidth}{l *{3}{>{\centering\arraybackslash}X} c}
        \toprule
        Конфигурация & Обуч. параметры, млн & RTF $\downarrow$ &
            Латентность, мс & Реальное время \\
        \midrule
        K1 (HuBERT+HiFi-GAN) & 18,8 & 0,031 & 310 & да \\
        K2 (HuBERT+WaveFit)  & 19,9 & 0,142 & 1420 & да \\
        K3 (WavLM+HiFi-GAN)  & 18,8 & 0,034 & 340 & да \\
        K4 (WavLM+WaveFit)   & 19,9 & 0,149 & 1490 & да \\
        \bottomrule
    \end{tabularx}
\end{table}

Конфигурации с HiFi-GAN работают примерно в $4{,}4$ раза быстрее аналогов с
WaveFit, что прямо следует из пятикратного итеративного вывода последнего.
При этом качество K3 (WavLM+HiFi-GAN) лишь незначительно уступает наиболее
тяжёлой K4. Сопоставление качества и быстродействия показано на
рис.~\ref{fig:tradeoff}.

\begin{figure}[H]
    \centering
    \begin{tikzpicture}
    \begin{axis}[
        width=0.82\textwidth, height=6.0cm,
        xlabel={RTF (меньше — быстрее)}, ylabel={PESQ},
        xmin=0, xmax=0.18, ymin=2.6, ymax=3.05,
        grid=both, grid style={black!12},
        nodes near coords, point meta=explicit symbolic,
        every node near coord/.append style={font=\scriptsize, anchor=south},
        tick label style={font=\scriptsize}, label style={font=\small}]
    \addplot[only marks, mark=*, mark size=2.2pt] coordinates {
        (0.031,2.71) [K1]
        (0.142,2.78) [K2]
        (0.034,2.94) [K3]
        (0.149,2.97) [K4]};
    \end{axis}
    \end{tikzpicture}
    \caption{Компромисс <<качество — быстродействие>> для четырёх конфигураций}
    \label{fig:tradeoff}
\end{figure}

\subsubsection{Сравнение с базовыми методами}

Чтобы оценить место предлагаемого подхода среди известных решений, рекомендуемая
конфигурация K3 сопоставлена с винеровской фильтрацией~\eqref{eq:wiener} и
дискриминативной моделью DEMUCS~\cite{demucs} на тестовом наборе B
(табл.~\ref{tab:baselines}). Классический фильтр обеспечивает умеренное
подавление шума, но не восстанавливает структуру сигнала; DEMUCS лучше по
SI-SDR, однако уступает по перцептивным метрикам PESQ и DNSMOS, поскольку не
выполняет ресинтез. Предлагаемая система превосходит оба базовых метода по
качеству восприятия.

\begin{table}[H]
    \caption{Сравнение с базовыми методами (тестовый набор B)}
    \label{tab:baselines}
    \centering
    \footnotesize
    \begin{tabularx}{\textwidth}{l *{4}{>{\centering\arraybackslash}X}}
        \toprule
        Метод & PESQ $\uparrow$ & STOI $\uparrow$ & SI-SDR $\uparrow$ &
            DNSMOS $\uparrow$ \\
        \midrule
        Вход (шум)            & 1,82 & 0,841 & 4,7  & 2,41 \\
        Винеровская фильтрация & 2,11 & 0,868 & 9,8  & 2,67 \\
        DEMUCS                & 2,64 & 0,915 & \textbf{16,1} & 3,12 \\
        K3 (наш подход)       & \textbf{2,94} & \textbf{0,927} & 15,2 & \textbf{3,41} \\
        \bottomrule
    \end{tabularx}
\end{table}

Качественное сравнение спектрограмм (схематично представлено на
рис.~\ref{fig:spectro}) показывает, что предлагаемая система не только подавляет
шумовой <<фон>> между гармониками, но и восстанавливает высокочастотную
гармоническую структуру, утраченную при сжатии и реверберации.

\begin{figure}[H]
    \centering
    \begin{tikzpicture}[font=\scriptsize]
        \def\W{3.6} \def\Ht{2.2} \def\G{0.6}
        % --- панель 1: чистый сигнал ---
        \def\xa{0}
        \draw[thick] (\xa,0) rectangle (\xa+\W,\Ht);
        \node[anchor=south] at (\xa+\W/2,\Ht) {Чистый $x$};
        \foreach \y in {0.35,0.75,1.15,1.55,1.95}
            \draw[line width=2.2pt, black!85] (\xa+0.1,\y) -- (\xa+\W-0.1,\y);
        % --- панель 2: искажённый сигнал (серый фон-шум + размытые гармоники) ---
        \def\xb{\W+\G}
        \begin{scope}
            \clip (\xb,0) rectangle (\xb+\W,\Ht);
            \fill[black!12] (\xb,0) rectangle (\xb+\W,\Ht);
            % видимая зернистость шума (детерминированная сетка)
            \foreach \i in {0,...,17} \foreach \j in {0,...,10}
                \fill[black!35] ({\xb+0.12+\i*0.2},{0.12+\j*0.2}) circle (0.7pt);
            % сохранились лишь нижние гармоники, размытые
            \foreach \y in {0.35,0.75,1.15}
                \draw[line width=2.6pt, black!45] (\xb+0.1,\y) -- (\xb+\W-0.1,\y);
        \end{scope}
        \draw[thick] (\xb,0) rectangle (\xb+\W,\Ht);
        \node[anchor=south] at (\xb+\W/2,\Ht) {Искажённый $y$};
        % --- панель 3: восстановленный сигнал ---
        \def\xc{2*\W+2*\G}
        \draw[thick] (\xc,0) rectangle (\xc+\W,\Ht);
        \node[anchor=south] at (\xc+\W/2,\Ht) {Восстановл. $\hat{x}$};
        \foreach \y in {0.35,0.75,1.15,1.55,1.95}
            \draw[line width=2.2pt, black!85] (\xc+0.1,\y) -- (\xc+\W-0.1,\y);
        % оси и стрелки
        \draw[->, thick] (\xa+\W,\Ht/2) -- (\xb,\Ht/2);
        \draw[->, thick] (\xb+\W,\Ht/2) -- (\xc,\Ht/2);
        \node[anchor=north] at (1.5*\W+\G,-0.2)
            {Частота (вертикаль) --- время (горизонталь)};
    \end{tikzpicture}
    \caption{Схематическое сравнение спектрограмм: восстановление подавленных
        гармоник и очистка межгармонического фона}
    \label{fig:spectro}
\end{figure}

\subsubsection{Анализ по типам деградаций}

Дополнительно конфигурация K3 проанализирована по типам искажений
(табл.~\ref{tab:ablation}). Наилучшее восстановление достигается при чисто
аддитивном шуме; добавление реверберации и кодека ожидаемо снижает метрики, но
система сохраняет существенный прирост во всех случаях.

\begin{table}[H]
    \caption{Качество конфигурации K3 по типам деградаций (PESQ / STOI)}
    \label{tab:ablation}
    \centering
    \footnotesize
    \begin{tabularx}{\textwidth}{X c c}
        \toprule
        Тип деградации & Вход & Выход K3 \\
        \midrule
        Шум                    & 1,95 / 0,86 & 3,08 / 0,94 \\
        Шум + реверберация     & 1,74 / 0,82 & 2,89 / 0,92 \\
        Шум + кодек            & 1,80 / 0,84 & 2,93 / 0,93 \\
        Шум + реверб. + кодек  & 1,61 / 0,79 & 2,79 / 0,90 \\
        \bottomrule
    \end{tabularx}
\end{table}

\subsection{Обсуждение результатов}

Эксперименты показывают, что закрытый энкодер USM может быть заменён открытыми
SSL-моделями без катастрофической потери качества. Решающим фактором оказался
выбор энкодера (WavLM предпочтительнее HuBERT) и слоя извлечения признаков, тогда
как тип вокодера влияет на качество слабее, но сильно --- на быстродействие.
Конфигурация \textbf{K3 (WavLM Base+ + HiFi-GAN, слой~7)} обеспечивает
оптимальный баланс: качество, близкое к наилучшему, при работе в десятки раз
быстрее реального времени и компактной обучаемой части. Полученное на открытых
данных DNSMOS ($3{,}41$--$3{,}52$) сопоставимо по порядку с приведёнными в
оригинальной работе~\cite{miipher2} значениями, что подтверждает
жизнеспособность открытой реализации.

\subsection{Ограничения}

Проведённое исследование имеет ряд ограничений. Во-первых, использованные
открытые энкодеры обучены преимущественно на англоязычной речи, поэтому
многоязычная универсальность, заявленная для USM, в данной реализации не
достигается. Во-вторых, оценка качества опиралась на объективные метрики;
субъективная экспертиза (MOS) не проводилась, а объективные метрики, как
известно, не полностью коррелируют с восприятием человека. В-третьих, обучение
велось на синтетических деградациях: поведение системы на реальных записях с
неизвестными типами помех требует отдельной проверки. Наконец, замороженный
энкодер ограничивает потенциал адаптации к специфике задачи --- совместное
дообучение могло бы дать дополнительный прирост качества ценой роста
вычислительных затрат.
